\section{The canonical comparison theorem}\label{sec:consequences}

Fix one computable admissible literal parsing policy \(\rho\). A
canonical assignment below means the zero-history assignment of the
literal input term under that policy, not the assignment of a normalized
term and not an arbitrary history-dependent assignment. Write
\(\can A_{f,\rho}\) for this vector.

\begin{theorem}[All ground free variables]\label{thm:ground}
For either \(f_*\) or a fixed \(f_b\), universal weak comparison,
Wilken--Weiermann strict comparison, and universal equality of canonical
assignments are decidable for original G\"odel T terms whose free variables
all have ground type. There is no bound on internal types, bound-variable
types, result types, recursor instances, or finite nesting.
\end{theorem}
\begin{proof}
The generated-source theorem gives, for the two supplied terms,
\[
 \can A_i=a_i,\quad\can B_i=b_i\quad(i>0),\qquad
 \can A_0=E_A(\mathbf x),\quad\can B_0=E_B(\mathbf x),
\]
where the positive coordinates are closed computable ordinals and the
scalar expressions are finite circuits built from positive natural inputs,
constants, sums, fixed natural multiples, and
\(u\mapsto\psi_f(\lambda+u)\) with fixed nonzero limit heads
\(\lambda<\varepsilon_0\). Repeated names have one shared coordinate.
Every positive integer is an admissible ground value, represented by a
finite sum of ones, so these inputs range independently over
\(\N_{>0}\).

First compute the finite positive-coordinate comparisons, using zero
padding beyond the assigned level. If any required one fails, the
vector comparison fails. Otherwise universal weak comparison is exactly
the absence of a solution to
\[
 \exists\mathbf x\in\N_{>0}^{\,r}\quad E_A(\mathbf x)>E_B(\mathbf x).
\]
For strict comparison the counterexample relation is \(\ge\), and
for equality it is \(\ne\). Flatten the scalar circuits using their
unique graph coordinates and true norm guards. Every graph coordinate
and every original input belongs to the same existential block.
Theorem~\ref{thm:arithmetic} decides this closed sentence. Internal
abstractions have already been handled by the finite assignment recursion;
they do not create external quantifier alternation. This proves all
three claims.
\end{proof}

The result is a forward translation theorem. It does not require a typed
source realizing every auxiliary interval created by arithmetic
elimination. Those intervals are legitimate arithmetic objects with the
explicit definitions proved above.

\begin{lemma}[Fixed-head tail growth]\label{lem:tail-growth}
For a nonzero limit \(\lambda\), let
\(h_\lambda(u)=\psi_f(\lambda+u)\). Then
\[
 h_\lambda(u+1)>2^{h_\lambda(u)}.
\]
\end{lemma}
\begin{proof}
Increasing the initial norm by one cannot shorten its nonfinite descent:
monotonicity of maximal predecessors, followed by taking limit parts,
propagates the weak comparison. Write the smaller terminal state as
\(Z=f^D(t)\). The larger terminal state is at least \(f^D(t+1)\).
At the last iterate, the preceding integer arguments differ by at least
one, so \eqref{eq:growth} gives \(f^D(t+1)\ge2^{2^Z}\).
Since \(D<Z\), the smaller collapse is \(Z+D<2Z\le2^Z\).
The terminal decomposition proves the claimed inequality.
\end{proof}

As an actual-source instance, with the native recursor argument order put
\(H=R^o0(D^o0):o\to o\). Its first positive coordinate is
\(\omega\cdot2\). Thus the application \(HX\) invokes the head
\(\omega^2\cdot2\) at the genuine raw norm
\(N(\omega^2\cdot2)+c_H+x_0\), where \(c_H\) is its actual closed
scalar. The complete-count table applies along this input ray without
replacing the source by an arbitrary ordinal-valued primitive.
